\begingroup
\fancyhead[L]{\small BEMERKUNGEN}
\renewcommand{\headrulewidth}{0pt}
\par\noindent\hspace*{\parindent}Zu [III.]--[IV.]: Die Rechnungen zu den speciellen Störungen finden sich
zum grössten Theil auf Blättern, zum kleinern in Handbüchern. Der gesamte Text
und zum Theil auch die Zahlen haben ergänzt werden müssen. Den Originalen sind
entnommen: S. 473--474 die Formeln, die sich auf den zur zweiten
\par\smallskip\noindent\hfill 76\textsuperscript{*}\par
\noindent\SourceLeaf{610}Rechnung der speciellen Störungen gehörigen Blättern finden,
S. 475--482 mit wenigen Ausnahmen die Zahlen, die theils auf Blättern, theils in
zwei Handbüchern (B\,c und B\,d) verstreut stehen und theilweise ergänzt wurden.
Die Zahlen der zweiten Rechnung, Seite 483--488 finden sich ohne Text,
aber zusammenhängend, auf einem Convolut von Blättern.

\par\noindent\hspace*{\parindent}Zu [V.]: Die Notizen [1.]--[3.] und [5.] haben vollständig frei bearbeitet
werden müssen, da im Nachlass gar keine Angabe über die Rechnungsmethode vorhanden
ist, sondern nur die Blätter mit den Rechnungen, nach deren Entzifferung erst die
Methode klargelegt werden konnte. Der Herausgeber wurde hiebei unterstützt durch
einzelne Vorarbeiten von \textsc{Schering} und \textsc{Schur}, welche bereits einige Blätter
entziffert und die Bedeutung einiger Zahlen angegeben hatten. Aus einigen
Aufzeichnungen \textsc{Scherings} liess sich auch erkennen, dass dieser sich vielleicht
den grössten Theil der Rechnungen klargelegt hatte; leider waren aber diese
Aufzeichnungen zu unvollständig, um bei der Bearbeitung benutzt werden zu können.
Das Rechnungsbeispiel [4.], das in einem Handbuch steht und mit nur unwesentlichen
Änderungen abgedruckt ist, diente als Anhaltspunkt. Die Bearbeitung ist so gehalten,
dass sie sich an die »Theoria interpolationis« Bd. III S. 265, anschliesst, da
dieses wahrscheinlich der von \textsc{Gauss} zunächst eingeschlagene Weg ist.
Die interessante Untersuchung in der »Quatrième Section« der »Exposition«, S. 465
kommt also hier garnicht zur Geltung. Ihr Zusammenhang mit der numerischen
Rechnung liegt aber auf der Hand, wenn die Integrationen zur Bestimmung der
Coefficienten der Reihen durch mechanische Quadratur ausgeführt werden.

\par\noindent\hspace*{\parindent}Von den numerischen Werthen der Notiz [6.] wurden die der Seiten 507--508
nach langem Suchen in einem Heftchen Schedae (A\,n) gefunden, die der Seiten 509--513
stehen auf einzelnen Blättern; sie mussten mehrfach ergänzt werden; in den Tafeln
für \(T\), \(V\), \(W\), S. 509--511 fehlen einige Werthe, auf deren Ergänzung verzichtet wurde.
Die Berechnung der Störungen selbst, Notiz [7.]--[12.], und die Vergleichung mit
den Oppositionen, Notiz [13.], stehen auf einzelnen Blättern. Es wurde nur soviel
abgedruckt, wie nöthig erschien, um eine lückenlose Darstellung zu erhalten.

\par\noindent\hspace*{\parindent}Auf einem der Blätter, welche die Integration der Störungen der Epoche
\(\varepsilon\) enthalten, findet sich folgende kleine Rechnung, wobei die erste der
angegebenen Zahlen die mittlere Bewegung der Pallas ist:
\begingroup
\newcommand{\BNoteJ}{\text{\mdseries\AstroSymbols ♃}}
\newcommand{\BNoteP}{\text{\mdseries\AstroSymbols ⚴}}
\[
\begin{gathered}
\mathit{769\overset{\prime\prime}{,}202079}\\[3pt]
\textit{das 7-fache}=\mathit{5384\overset{\prime\prime}{,}414553}\\[3pt]
\underline{\mathit{18}\ \textit{m.\,m.}\ \BNoteJ=\mathit{5384{,}392272}}\\[3pt]
\mathit{18}\BNoteJ-\mathit{7}\BNoteP=-\mathit{0{,}022281}\\[3pt]
\BNoteJ\ \textit{sid. in una die}\ \mathit{299\overset{\prime\prime}{,}132904}.
\end{gathered}
\]
\endgroup
\par\noindent Es scheint dies das Einzige zu sein, was auf den ersten Schritt der
\textsc{Gauss}schen Entdeckung der Rationalität beider Umlaufszeiten hinweist; daraus
lässt sich aber weiter nichts schliessen, als dass er eben die Grösse \(18n'-7n\)
als äusserst klein gefunden hatte.

\par\noindent\hspace*{\parindent}Die Werthe der Pallaselemente, die zum Grunde liegen, S. 508, hat
\textsc{Gauss} im Laufe der Rechnung zum Theil etwas verändert angenommen, so hat er
z. B. bei den Integrationen \(\tilde\omega-\AscendingNode=308^0\ 29'\ 23''\).
Einige kleinere Rechenfehler sind verbessert, eine Controlle der Rechnung ist
aber wegen der ungeheuren damit verbundenen Arbeit nicht vorgenommen worden.
\textsc{Gauss} scheint diese erste Rechnung der allgemeinen Störungen mit geringerer
Sorgfalt ausgeführt zu haben.
\clearpage
\endgroup
